Small students degrade when trained on reasoning models' long solutions, so recent work adjusts solution length with an interpolated teacher mixing a reasoning model and a short-answer instruct model at ratio $\lambda$. This teacher stops thinking altogether at some $\lambda$, but what this leaves in the student's weights is unknown. We ask what changes there, and whether that change carries the way of thinking. Across three teacher pairs and three students, we compare LoRA adapters trained at each $\lambda$, then remove, graft or amplify the region found and read outputs and SAE-based internal signals. Exactly where the teacher stops thinking, the student's adapter sharply changes the directions along which it updates the weights; before that, adapters always share the same small region. The region acts like a switch for thinking: removing it turns thinking off more than removing the same amount elsewhere does; amplifying it does not increase thinking; grafting it into a short-answer student starts thinking without finishing the solution. Removal also turns thinking off outside MATH, and much of the region is recovered in reasoning-distillation adapters trained without interpolation. Reasoning distillation into small students thus stores the teacher's way of thinking in a small, separate weight region that switches thinking on without carrying the solution through. More broadly, a change hidden in student outputs, where thinking and length collapse together, surfaces separately in LoRA-confined weights at the teacher's stopping $\lambda$.
